\section{Object-Storage-Native Search：耐久层与热查询层分离}

事故后的架构方向是把大块 immutable index segment 放入 object storage，把 query compute 和 SSD/memory cache 做成可替换层。Amazon S3 提供对象级耐久存储与强读后写一致性；turbopuffer 等系统在其上构建 search。优势是 durable cost 与 compute 解耦，代价是首次读取、tail latency、cache admission 和 segment management。

\subsection{Durable objects、cache 与 stateless compute}

\term{object storage} 以 object key 存放 blob 与 metadata，适合大对象、低成本耐久和海量 namespace，不提供传统 row update 与 join。\term{stateless compute} 指 query worker 不持有唯一耐久状态，可从 object 和 shared metadata 恢复；它仍可以有本地 cache，只是 cache 丢失不破坏 correctness。

\lecturefigure{10-object-storage-search.png}{Object-storage-native search 用 durable segment、SSD/memory cache 和 stateless query tier 分层。}{本地字幕 32:30--36:40；turbopuffer 官方架构资料。}

\begin{knowledgebox}{读图：Object storage 没有消灭 database 问题}
系统仍需要 metadata catalog、segment version、compaction、cache coherence、filter index、recall measurement 和 deletion。对象层降低 durability cost 并简化 replacement，却把性能问题移动到 prefetch、cache 和 query planning。
\end{knowledgebox}

\begin{importantbox}{Cache 的价值取决于 reuse}
若 object fetch latency 为 $L_o$，cache hit latency 为 $L_c$，hit rate 为 $h$，平均读取延迟近似
\[
E[L]=hL_c+(1-h)L_o.
\]
但用户体验更受 p95/p99 影响；热门 tenant、扫描 query 和 cold region 会让平均值掩盖尾部。需要按 namespace 与 query type 观察 hit rate。
\end{importantbox}

\begin{warningbox}{Low cost per TB 不等于 low query cost}
Object storage 的容量便宜，但 frequent GET、data transfer、decompression、CPU ranking 和 cache miss 会增加单位 query 成本。价格模型必须同时计算 stored bytes、read amplification、QPS、recall target 和 latency SLO。
\end{warningbox}

\subsection{本章小结}

Object-storage-native search 把 durable state 与 query compute 解耦，适合 immutable、可重建的大规模 index。它不是无数据库架构，而是重新定义 metadata、object 和 cache 的职责。

